%%%%%%%%%%%%%%%%%%%%%%%%%%%%%%%%%
%
%       EXERCÍCIO
%
%%%%%%%%%%%%%%%%%%%%%%%%%%%%%%%%%

\definevalorquestao

\begin{exercicioBanco}[\valorquestao]
O processo de resfriamento de uma liga metálica em um ambiente de fabricação é descrito por \(T(t) = T_{A} + \alpha \cdot 2^{\beta t}\), onde \(T(t)\) é a temperatura da liga metálica, em graus Celsius, no instante \(t\) (em minutos), \(T_{A}\) é a temperatura ambiente e \(\alpha\) e \(\beta\) são constantes. A referida liga metálica foi colocada em uma sala com temperatura ambiente de \(20^\circ\text{C}\). Um sensor acoplado ao metal indicou que ele atingiu \(148^\circ\text{C}\) após \(10\) minutos e chegou a \(52^\circ\text{C}\) após \(30\) minutos. Determine o valor de \(t\) para o qual a temperatura da liga metálica no recinto é apenas \(1^\circ\text{C}\) superior à temperatura ambiente.

\newcommand{\alternativas}{%
    \begin{center}
        \begin{tabularx}{\textwidth}{XXXXX}
            (a) \(50\) minutos. &
            (b) \(70\) minutos. &
            (c) \(80\) minutos. &
            (d) \(90\) minutos. &
            (e) \(120\) minutos.
        \end{tabularx}
    \end{center}
}

\newcommand{\resposta}{A}

% Lógica condicional para exibição
\ifdefstring{\atividade}{lista}{%
    \alternativas
    \vspace{0.5em}
    
    \noindent\textbf{Resposta:} letra \textbf{\resposta}.
}{%
    \ifdefstring{\modo}{objetivo}{%
        \alternativas
    }{%
        % modo = discursiva → não mostra alternativas
    }
}
\end{exercicioBanco}

